\section{Related Work and Open Problems}
\label{sec:related}

\subsection{Method of the Search}
\label{sec:search-method}

We did the search on 10 November 2026 with the Consensus search engine.
We used 15 queries. For each query we examined the first 20 records, thus 300 records in
total. \Cref{app:search} gives the queries. We read the title and the abstract of each
record. We kept a record if its abstract has one or more of the six properties that follow.

\begin{description}
  \item[A] One party makes decisions without a quorum during an outage.
  \item[S] The work has a succession rule.
  \item[H] The party with authority can be hostile.
  \item[L] The attacker controls the links, and thus the timeouts.
  \item[R] Removal tolerates hostile officers, and no message to the device is necessary.
  \item[T] The work gives a limit on the time that authority stays after removal.
\end{description}

We read parts of the full text of four sources only. Two are the paper on the Byzantine
generals problem~\cite{lamport1982byzantine} and the report on the 2015
attack~\cite{eisac2016ukraine}. The other two are the first section
of~\cite{babay2018spire} and pages 1 to 7 of~\cite{bommareddy2024compound}. For these
sources, the paper gives text from the source. The data about all other sources come from
their abstracts. \pending{Not done at this
time: read the full text of each source in \cref{tab:position}.}

\subsection{Five Groups of Work}

\paragraph{Byzantine agreement and replication}
The Byzantine generals problem has a commander who ``need not be
loyal''~\cite[p.~384]{lamport1982byzantine}. Its solutions make the loyal generals agree on
one order. They do not stop a hostile commander who gives the same dangerous order to all
generals. Replication protocols~\cite{castro1999pbft,castro2002pbft,amir2011prime} and
Byzantine quorum systems~\cite{malkhi1998byzquorum} give a leader no authority without a
quorum. A commander is different: the commander must make decisions without a quorum during
an outage.

Other work in this group is near the problem. One protocol puts a limit on the writes of a
hostile client after its removal~\cite{liskov2006byzclients}. Work on a change of the set of
servers makes a previous set safe after the change~\cite{kuznetsov2022reconfig}. Attacks on
timeouts are known. An attacker who delays messages causes incorrect timeouts and a change
of leader~\cite{tennage2023quepaxa,giuliari2024dos}. In replication, the new leader has no
authority without a quorum, thus the damage is to speed and not to safety.

\paragraph{Control systems for the grid that tolerate intrusion}
One group of designs replicates the SCADA server and tolerates hostile
copies~\cite{kirsch2014survivable,babay2018spire,babay2019deploying}. Subsequent work adds
substations~\cite{bommareddy2022realtime} and isolated sites~\cite{bommareddy2024compound}.
The operator stays loyal in this work. The authors of Spire write: ``Note that our threat
model does not cover a compromised HMI or a rogue operator''~\cite{babay2018spire}.

The work of 2024 gives the condition of the removal problem for servers. In automatic
removal, it records that ``compromised entities may try to exploit the reconfiguration
mechanism''~\cite{bommareddy2024compound}. For systems in which a quorum of servers removes
a server, it records that ``the voting criteria is left
undefined''~\cite{bommareddy2024compound}. This text is about servers. We found no
equivalent rule for the person who gives the orders.

\paragraph{Authentication of orders and access of operators}
Signatures show the sender of an
order~\cite{saxena2018signature,amoah2015dnp3,esiner2023provenance}. Other methods compare
the contents of an order with a simulation of the grid before the device obeys
it~\cite{meliopoulos2016command,mashima2018onthefly,mashima2019delaying,hong2019resilient}.
Access control for operators puts a limit on the substations that an operator
controls~\cite{nasr2018operator,nasr2024authorization}. Two designs of 2026 add a second
factor for an order~\cite{ahmad2026gridlogic} or a quorum of operators for a software
change~\cite{pawar2026certiflash}.

These methods examine one order: the sender and the contents. They do not examine if the
sender has command at this time, and they have no succession rule. A check of the contents
cannot reject an order that a loyal commander also gives (\cref{sec:limits}).

\paragraph{Leases, expiry and revocation}
A lease gives authority for a specified time~\cite{gray1989leases}. Three methods make sure that an access decision uses only new
data~\cite{stubblebine1995recent,krishnan2008stale,shakarami2019refresh}. Threshold methods
let a group of nodes give and revoke certificates in a network with no
server~\cite{zhou1999securing,luo2004ursa,liu2013ccrvc}. A survey gives the revocation methods for certificates on the
web~\cite{kwon2025revocation}. A second survey records that a check with no connection is
an open problem~\cite{lee2026ssirevocation}. One design gives a signed proof that a
credential was correct at a point in time~\cite{abraham2020revocable}.

HRCred is the nearest work that we found~\cite{zhou2026hrcred}. A threshold of parties
signs credentials with a short life for each epoch. Its authors show a maximum for the time
that a device accepts a credential after revocation. The credentials are for the identity
of a device. The abstract gives no succession rule and no minimum.

Some results show that two properties cannot be correct together, but they give no
quantity. The CAP result shows that a system
with a partition cannot have consistency and availability
together~\cite{gilbert2002cap,brewer2012cap}. An equivalent result is known for systems
that control equipment~\cite{lee2023cal}. Access control without a quorum makes decisions with
the data that a party has~\cite{jacob2025eventually,yanakieva2021impossibility}. Three
preprints examine the time at which authorization
stops~\cite{santosgrueiro2026effect,parakhin2026bureaucracy,deochake2026heartbeat}.

\paragraph{Handover between persons, and nuclear command}
Work on handover examines how control moves between operators, and between a person and an
automatic
system~\cite{flemisch2012balance,abbink2012haptic,grindley2024handover,nilsen2018handing,marot2022perspectives}.
All parties are loyal in the 20 records of that query. Nuclear command and control has the
two properties of this paper. An order must always go out when the correct authority gives it,
and it must not go out without that
authority~\cite{feaver1992guarding,crawford2025nc3}. A deputy with authority given before
the attack is one known solution~\cite{twigge1999learning}. This work gives the problem,
but it has no protocol model and no attacker on the links.

\begin{table*}[t]
\caption{Related work and the six properties of \cref{sec:search-method}. The entry ``part''
tells that the work has the property for a different party, or has only a part of it. We
made the entries from abstracts, and from text of the sources for the first and the third
row. A person must examine each entry against the full text.}
\label{tab:position}
\centering
\small
\begin{tabular}{@{}lcccccc@{}}
\toprule
Work & A & S & H & L & R & T \\
\midrule
Byzantine generals~\cite{lamport1982byzantine} & no & no & yes & no & no & no \\
Byzantine replication~\cite{castro1999pbft,castro2002pbft,amir2011prime} & no & yes & yes & part & part & no \\
Control systems that tolerate intrusion~\cite{babay2018spire,bommareddy2022realtime,bommareddy2024compound} & no & no & no & yes & no & no \\
Authentication and checks of orders~\cite{saxena2018signature,mashima2019delaying,hong2019resilient} & yes & no & part & no & no & no \\
Access control for operators~\cite{nasr2018operator,nasr2024authorization} & yes & no & yes & no & no & no \\
Leases~\cite{gray1989leases} & yes & part & no & no & part & maximum \\
Threshold certificates~\cite{zhou1999securing,luo2004ursa,liu2013ccrvc} & yes & no & yes & part & part & no \\
HRCred~\cite{zhou2026hrcred} & yes & no & yes & part & yes & maximum \\
Handover between operators~\cite{flemisch2012balance,grindley2024handover} & yes & yes & no & no & no & no \\
Nuclear command and control~\cite{feaver1992guarding,twigge1999learning} & yes & yes & part & no & no & no \\
\midrule
This paper & yes & yes & yes & yes & yes & minimum and maximum \\
\bottomrule
\end{tabular}
\end{table*}

\subsection{Open Problems}
\label{sec:gaps}

\Cref{tab:position} gives the result of the search. We found four open problems. For each,
we found no work in the 300 records and in the sources that we know. This is not a proof
that no such work is available.

\begin{description}
  \item[G1] \emph{Model.} No model has a commander who can be hostile and who must make decisions
  without a quorum during an outage. Replication has the hostile leader but not the
  decisions without a quorum. Leases have the decisions without a quorum but not the
  hostile commander.
  \item[G2] \emph{Limit.} No result gives a minimum for the tenure of a hostile commander. Leases
  and HRCred give a maximum. Thus a designer cannot know if a smaller tenure is possible.
  \item[G3] \emph{Timeout attack.} No succession rule puts a limit on the period that an attacker
  gets when the attacker causes a timeout. Work on replication examines incorrect timeouts
  only for speed.
  \item[G4] \emph{Removal.} No removal rule has these two properties together: it tolerates
  hostile officers, and no message to the device is necessary. For revocation by a quorum,
  a message is necessary. For expiry, no message is necessary, but known leases come from
  one server.
\end{description}

This paper gives a solution for each open problem. \Cref{sec:model} gives the model (G1).
\Cref{thm:limit} gives the minimum (G2). \Cref{thm:deputy} and rule D2 give the limit for
the timeout attack (G3). The rules of \sys{} and \cref{thm:tenure} give the removal rule
(G4).
